% ВВЕДЕНИЕ (2-5 страниц)
% Без жирных и курсивных шрифтов в тексте! (требование нормоконтроля)
\unnumberedsection{Введение}

В настоящее время процесс приёма абитуриентов в учреждения высшего образования Республики Беларусь претерпевает существенные изменения, связанные с переходом от бумажного документооборота к электронным системам обработки заявлений и автоматизированному расчёту результатов конкурса.
Современные приёмные комиссии оперируют тысячами заявлений, десятками конкурсных списков и сложными правилами распределения мест между льготными категориями, целевыми направлениями, бюджетной и платной формами обучения.
Ручная обработка такого объёма данных приводит к высоким трудозатратам сотрудников приёмной комиссии, увеличивает вероятность ошибок при подсчёте баллов и затрудняет оперативное получение сведений о текущем состоянии конкурса.

В Белорусском национальном техническом университете ежегодно подают заявления свыше двадцати тысяч абитуриентов, претендующих более чем на сто специальностей семнадцати факультетов.
Каждая специальность имеет собственный план приёма и набор конкурсных списков, разделённых на категории с фиксированными квотами и приоритетами вытеснения.
При этом один абитуриент может одновременно участвовать в нескольких конкурсах разных специальностей с указанием своего приоритета предпочтения.
Подобная многомерность задачи делает её формализуемой, но трудоёмкой при ручной обработке.

Актуальность дипломного проекта обусловлена необходимостью замены разрозненных электронных таблиц и устаревшего узкоспециализированного программного обеспечения единой централизованной системой, обеспечивающей: ведение справочников организационной структуры университета; регистрацию абитуриентов и их оценок по утверждённым критериям; ввод и редактирование заявлений с указанием приоритета предпочтения; автоматический пересчёт результатов отбора при любом изменении входных данных; разграничение прав доступа между операторами, аудиторами и администраторами; журналирование всех изменений данных для обеспечения прозрачности работы комиссии.

Цель дипломного проекта~--~разработка веб-приложения для управления приёмом абитуриентов в Белорусский национальный технический университет, обеспечивающего автоматизацию полного цикла работы приёмной комиссии от регистрации заявлений до формирования итоговых списков зачисленных.

Для достижения поставленной цели в проекте решаются следующие задачи:
\begin{enumerate}
\item проведён аналитический обзор существующих программных решений в области автоматизации работы приёмных комиссий и обоснован выбор средств реализации;
\item спроектирована клиент-серверная архитектура приложения, включающая одностраничное веб-приложение на React и серверную часть на ASP.NET~Core;
\item разработана реляционная модель данных с управлением версиями схемы через систему миграций Liquibase;
\item разработан и реализован алгоритм пересчёта результатов отбора с вытеснением слабых кандидатов в порядке защиты категорий и общего плана конкурсного списка;
\item реализована ролевая модель доступа на основе JSON Web Tokens с поддержкой пяти ролей пользователей;
\item спроектирована и реализована подсистема аудита, обеспечивающая журналирование всех изменений, валидацию данных вторым пользователем и двухступенчатое удаление с подтверждением;
\item реализован экспорт результатов отбора в формат Excel и поддержка двух языков интерфейса;
\item проведено функциональное тестирование разработанного приложения;
\item выполнено технико-экономическое обоснование разработки и рассмотрены вопросы охраны труда при работе с персональной электронно-вычислительной машиной.
\end{enumerate}

Краткое содержание расчётно-пояснительной записки.
В первом разделе проведён обзор предметной области, рассмотрены существующие аналоги программного обеспечения и сформулированы бизнес-, функциональные и технические требования к разрабатываемой системе.
Во втором разделе обоснован выбор стека технологий: фреймворка ASP.NET~Core для серверной части, библиотеки React совместно с компонентной системой Ant Design для клиентской части, системы управления базами данных PostgreSQL и инструмента миграций Liquibase.
В третьем разделе описаны архитектурные решения, реляционная модель данных, ролевая модель доступа и алгоритм пересчёта результатов отбора.
В четвёртом разделе приведены ключевые детали программной реализации обеих частей приложения, подсистемы аудита и двухступенчатого удаления.
Пятый раздел представляет собой руководство пользователя по работе с разработанным программным обеспечением.
В шестом разделе описана методика и результаты функционального тестирования.
Седьмой раздел содержит расчёты технико-экономического обоснования проекта.
Восьмой раздел посвящён вопросам охраны труда при работе с персональной электронно-вычислительной машиной.
